\documentclass[11pt]{article}
\usepackage[margin=1in]{geometry}
\usepackage[T1]{fontenc}
\usepackage{lmodern}
\usepackage{amsmath,amssymb,booktabs,graphicx}
\usepackage[colorlinks=true,allcolors=blue]{hyperref}
\newcommand{\StudyStatus}{\\[.5em]\large Working draft --- experiments incomplete}
\newcommand{\StudyAuthors}{Author details pending}
\newcommand{\StudyDate}{2026-09-25}

\title{Causal Interventions on Chain-of-Thought in Distilled and RL-Post-Trained Models\StudyStatus}
\author{\StudyAuthors}
\date{\StudyDate}
\begin{document}
\maketitle
\begin{abstract}
Visible reasoning and final answers can agree without the reasoning mediating the
model's decision. We study their relationship using two complementary interventions:
layer-specific activation steering and direct changes to generated reasoning text.
The protocol compares DeepSeek-R1-Distill-Qwen-32B with QwQ-32B, with smaller
distilled checkpoints as implementation pilots. It uses premise-disjoint SNLI contrast
pairs, per-example baselines, equal-norm random controls, complete generation records,
and explicit missing-output accounting. Reasoning assessment is blinded to final
answers and its consistency is reported separately from behavioral effects.
This working draft reports only completed implementation checks; the planned held-out checkpoint comparison remains incomplete.

The design measures intervention sensitivity and textual consistency; it does not
identify a pure effect of reinforcement learning or establish causal mediation from
reasoning--answer agreement alone.
\end{abstract}

\section{Introduction}
Chain-of-thought (CoT) is potentially useful for inspecting model behavior, but its
interpretation requires distinguishing plausible explanations from computations that
influence the answer. Behavioral interventions on reasoning reveal task-dependent
sensitivity \cite{lanham2023}, while naturally worded prompts can elicit internally
inconsistent justifications \cite{arcuschin2025}. Pre-CoT activation probes and
steering provide evidence that some answers can be influenced before verbal reasoning
begins \cite{cox2026}. An information-flow account makes the stronger distinction
between mediation through CoT and prompt-to-answer shortcuts \cite{jia2026}.

Our question is how hidden-state and text interventions differ across distilled and
RL-post-trained reasoning checkpoints. The intended contribution is a reproducible,
controlled behavioral comparison, with failure and missingness accounting. We avoid
equating an aligned answer change with faithful reasoning, or a robust answer with
post-hoc rationalization.

\section{Models and comparison limits}
The distilled checkpoints are DeepSeek-R1-Distill-Qwen-1.5B, 7B and 32B
\cite{deepseek2025}. The primary comparison is the 32B distilled checkpoint versus
the released QwQ-32B, whose official documentation lists both supervised fine-tuning
and reinforcement learning \cite{qwq2025}. Both use the Qwen2.5 base family at the
same parameter scale, but their post-training data and recipes differ. This is an
observational comparison of checkpoints under experimental inference-time
interventions, not random assignment to SFT versus RL.

Model revisions are pinned in the reproduction manifest. The planned 32B comparison
uses the same NF4 quantization with double quantization and bfloat16 computation.
The full-precision 1.5B run is an implementation pilot, not a quantization-controlled
scale comparison. Quantization can change the representation geometry and is a
limitation even when shared between models.

\section{Data and protocol}
We construct contrast pairs from SNLI \cite{bowman2015}. Each contains a premise,
an entailment hypothesis and a contradiction hypothesis. We sample 500 pairs from
the official training split, 50 from validation for calibration, and 50 from test
for evaluation. Premises normalized for case and whitespace are disjoint between
partitions. Candidate premises are shuffled before sampling. Invalid and neutral
labels are excluded. The task is a binary contrast subset, not the conventional
three-class SNLI benchmark. Dataset revision, seed and a checksum of all partitions
are recorded.

Prompts use each checkpoint's native chat template, an opening thought marker, and
the instruction to produce an explicit final label. Raw text, token IDs, stop reason,
prompt and seed are retained. Unfinished thought blocks, missing labels and invalid
formats remain separate from valid outcomes. Calibration precedes held-out evaluation;
all experimental settings and any changes must be recorded before test interpretation.

The held-out schedule was frozen locally after the 1.5B and 7B calibration runs;
it is not an external preregistration. It uses seed 42, temperature 0.6, top-$p$ 0.95,
top-$k$ 20 and batches of 16. The reasoning-generation ceiling is 2,048 new tokens,
increased from the calibration ceiling of 1,024 after observing incomplete traces.
Answer regeneration after text interventions is limited to 128 new tokens.
The same settings apply to the three held-out checkpoints. Every cell uses all 50
evaluation premises, with no replacement of examples whose baselines fail.

\subsection{Layer-specific activation steering}
For layer $l$, let $h_l(x)$ be the activation at the final prompt position before the
first generated reasoning token. Using only training pairs, compute
\begin{equation}
 d_l = \frac{1}{N}\sum_{i=1}^N h_l(x_i^{\rm contradiction})
       -\frac{1}{N}\sum_{i=1}^N h_l(x_i^{\rm entailment}).
\end{equation}
We add $\alpha d_l$ to the residual output of layer $l$ for the first $K=15$
prediction positions. Position one is the final prefill position, not the first
single-token decoding call. Activations are accumulated in float32. Intensity
is expressed in mean-difference units; it is not assumed to represent a universal
cross-model semantic dose.

We use intensities $\alpha\in\{0.5,1,2,4\}$ and relative layer positions
$\{0,.15,.30,.45,.60,.75,.90,1\}$, rounded to the nearest zero-based layer index.
The 28-layer checkpoint therefore uses layers $\{0,4,8,12,16,20,24,27\}$; the
64-layer checkpoints use $\{0,9,19,28,38,47,57,63\}$. Each model has 32 semantic
steering cells and 64 matched random-control cells. Each cell contains one seeded
generation per premise; the study does not estimate generation-seed variability.

Every evaluated example has an unsteered baseline. For each layer and intensity,
controls use (i) a Gaussian direction projected orthogonally to $d_l$ and normalized
to the same norm and (ii) a Gaussian direction normalized to the same norm without
projection. Each control direction is fixed across its 15 positions and separately
seeded by example and layer. These controls distinguish a directional effect from
perturbation magnitude but cannot eliminate every off-manifold effect.

\subsection{Direct interventions on written reasoning}
For completed baseline traces, retain fractions $q\in\{0,.25,.50,.75,1\}$ of
reasoning tokens, close the thought block and regenerate an answer. The $q=1$
condition controls for the answer-regeneration procedure. We also substitute the
reasoning from the paired contradiction-hypothesis prompt and separately shuffle the
original reasoning tokens. These operations test answer sensitivity to the reasoning
prefix. They do not preserve all distributional properties, and the substituted
trace may contain information inconsistent with the original hypothesis.

\subsection{Reasoning assessment}
A separate model judges the conclusion supported by the reasoning text, without
seeing the final answer. Its outputs are restricted to entailment, contradiction,
mixed or unclear. Exact supporting quotations and JSON fields are validated; errors
are preserved. A quote-based first-contradiction token index is an approximate textual
indicator, not internal commitment time. Synthetic cases check extraction, negation,
self-correction and empty inputs. This validation is not a substitute for independent
human annotation of model-generated traces. Retest agreement is computed on a seeded
sample; undefined Cohen's kappa is reported as unavailable.

We derive descriptive reasoning--answer categories only after assessment. An
entailment baseline that changes to a contradiction answer can either have aligned
contradiction reasoning or a reasoning--answer mismatch. These labels are not proofs
of faithfulness or deception.

\section{Analysis}
Report baseline accuracy and completion separately. Among examples whose baseline
answer is entailment, report the contradiction-answer rate and the number of missing
intervention answers. If $f$ observed flips and $m$ missing answers occur among $n$
eligible examples, the missingness bounds are $[f/n,(f+m)/n]$. Complete-case rates
receive Wilson intervals and are explicitly subject to selection bias when missingness
depends on intervention strength.

Control differences are paired by example and bootstrapped over examples, using
5,000 resamples. Layers and intensity settings on the same example are not independent
samples. These are exploratory intervals; the many cells require caution about
multiplicity and selected extremes. The final checkpoint comparison must use shared
examples, aligned relative layer positions and identical operating settings.

Supplementary analyses compare each text intervention with the regenerated answer
after the full original prefix, and quantify the full-prefix regeneration effect
itself. These comparisons use both all eligible baselines and the originally correct
baseline subset. Balanced baseline accuracy resamples the two hypotheses sharing
a premise as one cluster. Activation/control comparisons additionally receive
two-sided exact tests based on discordant pairs, with Holm adjustment within each
model's activation/control family. Complete-case selection remains a limitation;
missing-outcome bounds use the observed label in a partially missing pair where
available. Cross-model and text comparisons remain exploratory.

\section{Execution status and results}
\noindent\textbf{Draft status:} The held-out comparison is not yet complete. The table below is an execution audit, not evidence of a training-method effect.
\begin{center}
\begin{tabular}{lrrrr}
\toprule
Run & Examples & Valid & Correct & Generations \\
\midrule
distill-1.5b / calibration & 4 & 1 & 0 & 51 \\
distill-1.5b / calibration & 20 & 16 & 6 & 40 \\
distill-7b / calibration & 20 & 16 & 16 & 40 \\
\bottomrule
\end{tabular}
\end{center}
Valid denotes a parseable final answer. Correct uses the dataset label and counts missing answers as incorrect. Calibration observations are not held-out findings.
The assessor passed 12 of 12 synthetic validation cases. These are engineered checks, not human judgments of generated traces.


\section{Limitations}
This study concerns a small binary NLI subset and cannot establish general-purpose
reasoning monitorability. SNLI labels can encode pragmatic assumptions, and public
benchmark data may have appeared in pretraining. Baseline task failure limits
interpretation of attempted answer shifts. Limited generation budgets can introduce
condition-dependent missingness. Model-family similarities between an assessor and
evaluated checkpoints can introduce correlated judgments. Synthetic validation and
retest agreement do not establish external validity.

Hidden-state interventions can affect reasoning and answers through separate paths.
Consequently, coordinated changes do not identify mediation. An unchanged answer
after truncation may reflect redundant reasoning; a changed answer may reflect
an unnatural prefix. The checkpoint comparison also confounds training recipe,
training data and other post-training choices. We therefore do not infer a universal
parameter threshold, an irrevocable commitment layer, or a causal advantage of
SFT, RL, outcome rewards or process rewards.

\section{Reproducibility}
The repository supplies dataset and model revisions, per-run configuration hashes,
source hashes, software versions, raw generations, preserved assessor outputs,
analysis code and tests. Resume rejects a changed run identity. Generation hooks are
removed even on failures. The historical v3 outputs are excluded because their
reasoning extraction was empty and their evaluation reused one example. Source and
result checks must pass before a submission archive can be considered complete.

\bibliographystyle{plain}
\bibliography{references}
\end{document}
